In this section we will report some results obtained by training the neural
model and applying it to an MPC based trajectory tracking problem. 

\subsection{Neural Network Training}

\subsubsection{Derivative Prediction}

In figure \ref{fig:derivative_loss} we can observe that the loss descends
pretty quickly, which means that the network is converging quickly.

\begin{figure}[H]
    \centering
    \missingfigure{Loss plot for the derivative prediction task}
    % \includegraphics[width=0.8\textwidth]{media/derivative_loss.png}
    \caption{Loss plot for the derivative prediction task}
    \label{fig:derivative_loss}
\end{figure}

This reflects also in a successful evaluation on the test set, where
predictions are very close to the nominal model, with very few exceptions.
Nonetheless, the performance inside the MPC is not as good as expected. As
an example, we report here a figure depicting how the unicycle drifts away
from its desired path (a straight line). This can be mitigated only by
putting very high weights on the control effort and tahe position error,
which is surely not an ideal situation, also because the nominal MPC works
fine with both sets of weights. This bad performance could be imputable to
an overfgitting situation: since the simulation is inherently noisy, since
we need to discretize everything, maybe the network is not able to grasp
that, being trained just on the current state, without any notion on
discretization timesteps or numerical integration methods.

\begin{figure}[H]
    \centering
    \missingfigure{Drift of the unicycle from the desired path}
    % \includegraphics[width=0.8\textwidth]{media/drifting.png}
    \caption{Drift of the unicycle from the desired path}
    \label{fig:drifting}
\end{figure}


\subsubsection{State Prediction}

The loss for the state prediction task is depicted in figure \ref{fig:state_loss}.

\begin{figure}[H]
    \centering
    \missingfigure{Loss plot for the state prediction task}
    % \includegraphics[width=0.8\textwidth]{media/state_loss.png}
    \caption{Loss plot for the state prediction task}
    \label{fig:state_loss}
\end{figure}




